Insee, comptes nationaux 2025 (tableau 8.212) ; DGFiP Statistiques n° 21 ; CPO (2025) ; calculs de l'auteur (fiscalite/is/simulation\_is.py).\par Note : coût annuel pour l'État, hors intérêts, d'une déduction immédiate des machines, équipements et matériels informatiques des sociétés non financières (84 Md\euro{} en 2025, hors matériels de transport), à partir de 2027, croissance nominale de 3 \% par an, 75 \% de cet investissement fait par des sociétés imposables (hypothèse). Bande de chaque coût : régime actuel supposé dégressif (bas) ou linéaire sur sept ans (haut). Sans investissement supplémentaire suscité par la mesure. Ressources, en euros de 2027 puis en hausse de 3 \% par an : élargissement de l'assiette, 3,5 à 4,9 Md\euro{} (fin du taux réduit de 15 \%, CIR apprécié au niveau du groupe, fin des exonérations zonées hors outre-mer, avantage de la boîte à brevets réduit de moitié) ; taux à 28 \%, 5,1 à 8,7 Md\euro{} selon la réaction des bénéfices déclarés (0 à 1,5 \% de base en moins par point de taux). Le taux réduit et la hausse de taux sont chiffrés sur les bénéfices de 2022, portés en 2027 par la croissance nominale.
